\documentclass[12pt,a4paper]{article}
\usepackage[utf8]{inputenc}
\usepackage[english]{babel}
\usepackage[T1]{fontenc}
\usepackage{amsmath, amssymb, amsfonts}
\usepackage{graphicx}
\usepackage{booktabs}
\usepackage{hyperref}
\usepackage{geometry}
\geometry{margin=2.5cm}
\usepackage{caption}
\usepackage{float}
\usepackage{siunitx}

\title{Decoupled Observer Patch Holography (dOPH): \\
Mathematical Foundation and First Astrophysical Tests}
\author{v04226079}
\date{\today}

\begin{document}

\maketitle

\begin{abstract}
We present the Decoupled Observer Patch Holography (dOPH) framework --- a scale-dependent modified gravity theory based on holographic principles and dynamic decoupling of observer patches. 

A central result is the derivation of the universal scaling factor 
\[
P = \sqrt{\frac{8}{3}} \approx 1.633
\]
as a geometric fixed point emerging from the minimization of a mismatch functional on patch interfaces. 

The model is tested on the massive elliptical galaxy NGC 4649 (M60). Using only a realistic stellar mass-to-light ratio (\( M/L \approx 8 \)), dOPH successfully reproduces the nearly flat globular cluster velocity dispersion profile up to \( \sim 38 \) kpc. The theory performs competitively with MOND and dark matter halo models while requiring significantly fewer free parameters. Strengths, limitations, and future directions (especially regarding the Bullet Cluster) are discussed.
\end{abstract}

\section{Introduction}

The observed dynamics of galaxies and clusters continue to challenge our understanding of gravity. While $\Lambda$CDM remains successful on large scales, modified gravity theories offer alternative explanations.

This paper develops the \textbf{Decoupled Observer Patch Holography (dOPH)} framework, in which spacetime consists of finite observer patches with holographic screens and dynamic decoupling mechanisms.

\section{Mathematical Derivation of \( P = \sqrt{8/3} \)}

Consider the interface between two adjacent observer patches with natural accelerations \( g_1 \) and \( g_2 \). We minimize the mismatch functional:

\begin{equation}
J(\lambda) = (g_1 - \lambda g_2)^2 + \beta \left( \lambda g_2 - \frac{g_1}{\lambda} \right)^2
\end{equation}

In the approximation \( g_1 \approx g_2 = g \) this reduces to:

\begin{equation}
f(\lambda) = (1 - \lambda)^2 + \beta \left( \lambda - \frac{1}{\lambda} \right)^2
\end{equation}

With the geometrically motivated coefficient \( \beta = 3/8 \), the minimum occurs exactly at 

\[
\lambda = P = \sqrt{\frac{8}{3}} \approx 1.633.
\]

This value is a **geometric fixed point** arising from optimal 3D packing (HCP/FCC), interface mismatch minimization, and conformal considerations.

\section{Test on NGC 4649 (M60)}

NGC 4649 shows a nearly flat GC velocity dispersion \( \sigma_{\rm LOS} \approx 220 \)--$235$ km/s up to \( \sim 38 \) kpc. Using Hernquist stellar profile with \( M/L \approx 8 \) and spherical Jeans equation, dOPH successfully reproduces the observed profile thanks to scale-dependent enhancement. The model outperforms pure Newtonian gravity and is competitive with MOND and NFW models with far fewer parameters.

\section{Discussion and Future Work}

dOPH performs well on galactic scales but requires improvements for cluster collisions (e.g., Bullet Cluster). Future work includes non-local terms, multi-component decoupling, and time-dependent dynamics.

\section{Conclusion}

We derived \( P = \sqrt{8/3} \) as a geometric fixed point and demonstrated promising performance on NGC 4649. dOPH offers a compelling holographic, scale-dependent approach to gravity.

All code is available at: \url{https://github.com/v04226079-sketch/dOPH-Decoupled-Observer-Patch-Holography}

\end{document}
